
le nombre des polynômes @@M0@@ est au plus égal au nombre de composantes connexes de @@M1@@ (EGA III 7.9.5) et c'est un petit lemme que si pour un @@M2@@, @@M3@@ est localement libre de rang @@M4@@, alors @@M5@@ l'est aussi au-dessus d'un voisinage ouvert de @@M6@@ dans @@M7@@, (voir EGA IV 11.3.10, où on fait @@M8@@, @@M9@@). Or, il résulte des hypothèses et de (EGA III 7.7.6) que quitte à couper @@M10@@ encore plus fin, on peut trouver une suite @@M11@@ de sections @@M12@@ qui est @@M13@@-régulière. Donc (ii) résulte de la propriété de semi-continuité supérieure de la fonction @@M14@@ où pour @@M15@@ et @@M16@@, @@M17@@ désigne la restriction de @@M18@@ au schéma des zéros de @@M19@@ (EGA III 7.7.5).

L'implication (ii) @@M20@@ (iii) résulte immédiatement de 1.11. L'implication (iii) @@M21@@ (iv) résulte de 1.3 (iii) en prenant @@M22@@ et @@M23@@.

Supposons que @@M24@@ satisfasse à (iv); il existe des suites exactes
@@M25@@
et l'ensemble des polynômes @@M26@@ est fini. Par hypothèse la famille des classes des @@M27@@ est limitée; donc par l'implication (i) @@M28@@ (ii) déjà prouvée, l'ensemble des polynômes @@M29@@ est fini et il existe une suite d'entiers @@M30@@ telle que tous les @@M31@@ (avec @@M32@@ algébriquement clos) soient des @@M33@@-faisceaux. Par suite, l'ensemble des polynômes @@M34@@ est fini et en vertu de 1.6, tous les @@M35@@ sont des @@M36@@-faisceaux. Donc en appliquant à la famille des classes des @@M37@@ l'implication (ii) @@M38@@ (iv) déjà prouvée, on trouve (v).

Supposons enfin que @@M39@@ satisfasse à (v), prouvons que @@M40@@ est limitée. On peut supposer @@M41@@ (resp. @@M42@@) de la forme @@M43@@. En effet,
